% Chapter 2: Supporting Results
\chapter{Supporting Results}

This chapter collects two auxiliary results that are used in later chapters but do not
fit naturally into the main proof chain.  The first is a spectral containment lemma for
conjugate-weighted sums, which is fully proved in Lean.  The second is a commutativity
identity for real powers in a C$^*$-algebra, which is currently admitted pending
upstream Mathlib infrastructure.

\section{Conjugate-Weighted Sums}

\begin{lemma}[Spectrum of a conjugate-weighted sum]
  \label{lem:conj_sum_spectrum}
  \lean{spectrum_sum_conj_subset}
  \leanok
  Let $\mathcal{B}$ be an ordered unital C$^*$-algebra, let $I \subseteq \mathbb{R}$
  be a convex subset, and let $n \in \mathbb{N}$.  Suppose
  $a_1, \dots, a_n \in \mathcal{B}$ are self-adjoint with $\sigma(a_i) \subseteq I$
  for each $i$, and let $b_1, \dots, b_n \in \mathcal{B}$ satisfy
  $\sum_i b_i^* b_i = 1$.  Then
  \[
    \sigma\!\left(\sum_i b_i^*\, a_i\, b_i\right) \subseteq I.
  \]
  \begin{proof}
    Set $\alpha = \inf_i \inf \sigma(a_i)$ and $\beta = \sup_i \sup \sigma(a_i)$;
    then $[\alpha, \beta] \subseteq I$ by convexity and ord-connectedness.
    The unit constraint $\sum_i b_i^* b_i = 1$ gives
    $\alpha \cdot 1 \leq \sum_i b_i^* a_i b_i \leq \beta \cdot 1$
    in the order on $\mathcal{B}$, since each summand satisfies
    $\alpha \cdot b_i^* b_i \leq b_i^* a_i b_i \leq \beta \cdot b_i^* b_i$.
    The order bounds imply the spectrum is contained in $[\alpha, \beta]$,
    and hence in $I$.
  \end{proof}
\end{lemma}

\section{Real Powers in C\texorpdfstring{$^*$}{*}-Algebras}

\begin{lemma}[Power of a product of commuting positives]
  \label{lem:mul_rpow_commute}
  \lean{CFC.mul_rpow_of_commute}
  Let $\mathcal{B}$ be an ordered unital C$^*$-algebra, let $a, b \in \mathcal{B}$
  with $0 \leq a$, $0 \leq b$, and $ab = ba$.  Then for all $r \in \mathbb{R}$,
  \[
    (ab)^r = a^r\, b^r,
  \]
  where the powers are taken via the continuous functional calculus.
  \begin{proof}
    \emph{Admitted.}  This identity would follow from the fact that the continuous
    functional calculus distributes over multiplication for commuting elements.
    Specifically, one would need that $\mathrm{cfc}(f, ab) = \mathrm{cfc}(f, a)\,
    \mathrm{cfc}(f, b)$ whenever $a$ and $b$ commute and $f$ is the function
    $t \mapsto t^r$.  This CFC multiplicativity for commuting pairs is not yet
    available in Mathlib, so the lemma is currently admitted as an axiom.
  \end{proof}
\end{lemma}
